% Report-ready section for the NP-completeness component.
\section{Cloud Workload Consolidation and Bin Packing}
\label{sec:binpacking}

\subsection{Problem abstraction}
Model each workload as an item with positive size $a_i$ and each identical
server as a bin of capacity $B$. The objective is to assign every item to
a bin without exceeding capacity and to minimize the number of bins used.
The decision version asks whether all items can be packed into at most $k$
bins. This is a simplified one-dimensional consolidation model; it does
not capture multi-dimensional resources, placement constraints, or runtime
migration costs.

\subsection{Membership in NP}
A certificate specifies the bin assigned to each item. In polynomial time,
a verifier checks that no item is assigned outside the proposed bins and
that the sum of item sizes in each bin does not exceed $B$. Therefore,
the decision version of Bin Packing is in NP.

\subsection{NP-hardness reduction from 3-Partition}
The complexity classification and reduction framework
follow the standard theory of NP-completeness
\cite{gareyjohnson}.
We reduce from 3-Partition, a strongly NP-complete problem. An instance
contains $3m$ positive integers $a_1,\ldots,a_{3m}$ and a positive
integer $B$ satisfying $\sum_i a_i=mB$ and
$B/4<a_i<B/2$ for every item. The question is whether the integers can
be partitioned into $m$ triples, each summing to $B$.

Construct a Bin Packing decision instance using the same $3m$ item sizes,
bin capacity $B$, and target number of bins $k=m$. This construction is
polynomial in the input size.

\paragraph{Forward direction.}
If the 3-Partition instance has a partition into $m$ triples each summing
to $B$, place each triple in one bin. Every bin has capacity $B$ and
exactly $m$ bins are used. Thus the constructed Bin Packing instance is
a yes-instance.

\paragraph{Reverse direction.}
Suppose the constructed Bin Packing instance can pack the items into at
most $m$ bins. The total item size is $mB$, while each bin has capacity
$B$. Consequently, all $m$ bins must be used and each must be filled to
capacity $B$. Since every item is strictly larger than $B/4$ and
strictly smaller than $B/2$, no bin can contain four or more items and
no bin can contain two or fewer items while being full. Each bin must
therefore contain exactly three items summing to $B$. These triples form
a valid solution to the 3-Partition instance.

\paragraph{Conclusion.}
The reduction preserves yes- and no-instances, so Bin Packing decision is
NP-hard. Combined with membership in NP, it is NP-complete. Because the
reduction is from strongly NP-complete 3-Partition, Bin Packing is also
strongly NP-hard in its optimization form.

\subsection{Greedy heuristic: First Fit Decreasing}
First Fit Decreasing is a classical bin-packing heuristic
studied in the literature on approximation algorithms
\cite{johnson1974}.
The implemented heuristic sorts items in nonincreasing order and places
each item into the earliest opened bin with sufficient remaining capacity.
If no opened bin fits, it opens a new bin. A segment tree stores the
maximum remaining capacity in each subtree and finds the leftmost bin
whose remaining capacity is at least the current item's size; the tree is
updated after each placement. The leftmost-fit query preserves First Fit
behavior while accelerating the search.

\subsection{Running time}
Sorting takes $O(n\log n)$. Each item requires a segment-tree query and
update, each taking $O(\log n)$, so placement takes $O(n\log n)$ total.
The overall running time is $O(n\log n)$ and the storage is $O(n)$.
First Fit Decreasing is a heuristic and does not guarantee an optimal
packing on every instance.